\begin{thebibliography}{99}

\bibitem{guinand}
Guinand, A. P., 
\textit{The umbral method: a survey of elementary mnemonic and manipulative uses}, 
Amer. Math. Monthly 86 (1979), no. 3, 187–195.
\url{https://doi.org/10.1080/00029890.1979.11994763}

\bibitem{bucchianico}
Di Bucchianico, A.; Loeb, D., 
\textit{A selected survey of umbral calculus}, Electron. J. Combin. 1 (1995), Dynamic Survey 3, 34 pp.

\bibitem{roman1}
Roman, Steven, \textit{The umbral calculus}, 
Pure and Applied Mathematics, 111, Academic Press, Inc., London, 1984.

\bibitem{roman2}
Roman, Steven M.; Rota, Gian-Carlo, \textit{The umbral calculus}, 
Adv. in Math. 27 (1978), no. 2, 95–188.
\url{https://doi.org/10.1016/0001-8708(78)90087-7}

\bibitem{rota}
Rota, G.-C.; Taylor, B. D., \textit{The classical umbral calculus}, 
SIAM J. Math. Anal. 25 (1994), no. 2, 694–711.

\bibitem{bell}
Bell, E. T., \textit{The history of Blissard’s symbolic method, with a sketch of its inventor’s life}, 
Amer. Math. Monthly 45 (1938), no. 7, 414–421.
\url{https://doi.org/10.1080/00029890.1938.11990829}

\bibitem{kim1}
Kim, Dae San; Kim, Taekyun, \textit{Umbral calculus associated with Bernoulli polynomials}, J. Number Theory 147 (2015), 871–882.
\url{https://doi.org/10.1016/j.jnt.2013.09.013}

\bibitem{blissard}
Blissard, John, \textit{Theory of generic equations}, 
Quart. J. Pure Appl. Math. 5 (1862), 58–75, 185–208.

\bibitem{verdoodt}
Verdoodt, A., \textit{p-adic q-umbral calculus}, 
J. Math. Anal. Appl. 198 (1996), no. 1, 166–177.

\bibitem{sambale}
Sambale, B., \textit{An invitation to formal power series}, 
Jahresber. Dtsch. Math.-Ver. 125 (2023), no. 2, 69–110.
\url{https://doi.org/10.1365/s13291-022-00256-6}

\bibitem{kim2}
Kim, D., \textit{A class of Sheffer sequences of some complex polynomials and their degenerate types}, Mathematics 7 (2019), no. 11, Paper No. 1064, 13 pp.
\url{https://doi.org/10.3390/math7111064}

\bibitem{niederhausen}
Niederhausen, H., \textit{Rota’s umbral calculus and recursions}, 
Algebra Universalis 49 (2003), no. 4, 435–457.
\url{https://doi.org/10.1007/s00012-003-1820-6}

\bibitem{costabile}
Costabile, Francesco, \textit{Modern umbral calculus: an elementary introduction with applications to linear interpolation and operator approximation theory}, De Gruyter, Berlin, 2019.
\url{https://doi.org/10.1515/9783110652925} 

\bibitem{davis}
Davis P. J., \textit{Interpolation and approximation}, Dover (1975).

\bibitem{roman3}
Roman, Steven,
\textit{ More on the umbral calculus, with emphasis on the q-umbral calculus}, J. Math. Anal. Appl. 107 (1985), no. 1, 222–254.
\url{https://doi.org/10.1016/0022-247X(85)90367-1.}

\bibitem{roman4}
Roman, Steven, \textit{The theory of the umbral calculus. I}, 
J. Math. Anal. Appl. 87 (1982), no. 1, 58–115.
\url{https://doi.org/10.1016/0022-247X(82)90154-8.}

\bibitem{razpet}
Razpet, Marko, \textit{An application of the umbral calculus}, 
J. Math. Anal. Appl. 149 (1990), no. 1, 1–16.
\url{https://doi.org/10.1016/0022-247X(90)90281-J}

\bibitem{gessel}
Gessel, I. M., \textit{Applications of the classical umbral calculus}, Algebra Universalis 49 (2003), no. 4, 397–434.
\url{https://doi.org/10.1007/s00012-003-1813-5}

\end{thebibliography}